% 等比例七站覆盖图：坐标单位 km，标注单位 m，整体旋转不影响覆盖。
\begin{tikzpicture}[x=1.5cm,y=1.5cm,font=\small,>=Latex]
\definecolor{p3coverblue}{RGB}{46,111,164}
\definecolor{p3coverred}{RGB}{176,55,45}
\coordinate (S0) at (0,0);
\foreach \j in {1,...,6}{
  \pgfmathsetmacro{\ang}{(\j-1)*60}
  \coordinate (S\j) at (\ang:1.14);
}
% 先绘填充，再绘圆周，避免后画的圆遮盖先画的线。
\foreach \j in {0,...,6}{\fill[p3coverblue,opacity=0.09] (S\j) circle[radius=1];}
\foreach \j in {0,...,6}{\draw[p3coverblue,line width=0.45pt] (S\j) circle[radius=1];}
\draw[black,line width=1pt] (0,0) circle[radius=1.8];
\draw[black!50,densely dashed,line width=0.45pt] (0,0) circle[radius=1.14];
\foreach \j/\pos in {1/below right,2/above,3/above,4/below left,5/below,6/below}{
  \fill[p3coverblue] (S\j) circle[radius=0.033];
  \node[\pos,inner sep=2pt] at (S\j) {$S_{\j}$};
}
\fill[p3coverblue] (S0) circle[radius=0.035];
\node[above left,inner sep=2pt] at (S0) {$S_0$};
% 半径与最不利边界位置。
\draw[->,black!65,line width=0.45pt] (0,0)--(235:1.8)
  node[pos=0.75,sloped,above,fill=white,inner sep=1pt,font=\footnotesize] {$1800\,\mathrm m$};
\draw[<->,black!65,line width=0.45pt] (0,0)--(S1)
  node[midway,below,fill=white,inner sep=1pt,font=\footnotesize] {$1140\,\mathrm m$};
\coordinate (Gstar) at (30:1.8);
\draw[p3coverred,dashed,line width=0.8pt] (S1)--(Gstar)--(S2);
\fill[p3coverred] (Gstar) circle[radius=0.04];
\node[above right,inner sep=2pt] at (Gstar) {$G_*$};
% 图右侧短图例，不压住圆域。
\draw[black,line width=1pt] (2.45,1.6)--(2.85,1.6);
\node[anchor=west] at (2.96,1.6) {场地边界：$1800$ m};
\draw[p3coverblue,line width=0.65pt] (2.45,1.14)--(2.85,1.14);
\node[anchor=west] at (2.96,1.14) {各站接收范围：$1000$ m};
\draw[black!50,densely dashed] (2.45,0.68)--(2.85,0.68);
\node[anchor=west] at (2.96,0.68) {外圈布设半径：$1140$ m};
\fill[p3coverblue] (2.65,0.22) circle[radius=0.035];
\node[anchor=west] at (2.96,0.22) {检测点 $S_0,\ldots,S_6$};
\node[anchor=north west,align=left,inner sep=0pt] at (2.45,-0.36)
  {最不利点示例：$G_*$\\[3pt]
   $\|G_*-S_1\|=\|G_*-S_2\|$\\[3pt]
   $\phantom{\|G_*-S_1\|}\approx992.689\ \mathrm m<1000\ \mathrm m$};
\node[anchor=north west,align=left,inner sep=0pt] at (2.45,-1.42)
  {七个接收圆的并集覆盖整个场地\\[3pt]等比例绘制，取 $\phi=0^\circ$};
\end{tikzpicture}
